\section{Recommended Reading}

\begin{readingbox}
\textbf{Foundation}
\begin{itemize}[nosep]
  \item John Holland, \emph{Emergence}, for how repeated local interaction can
        consolidate higher-level patterns.
  \item Eugene Izhikevich, \emph{Dynamical Systems in Neuroscience}, for
        attractor, basin, and transition language.
  \item Stanislas Dehaene, \emph{Consciousness and the Brain}, for a readable
        bridge from state-based neuroscience to practical interpretation.
\end{itemize}
\textbf{Modern Research}
\begin{itemize}[nosep]
  \item \emph{Metastability demystified: the foundational past, the pragmatic
        present and the promising future} (Nature Reviews Neuroscience, 2024,
        medium).
  \item \emph{Associative conditioning in gene regulatory network models
        increases integrative causal emergence} (Communications Biology, 2025,
        medium).
  \item \emph{Cortical connectivity, local dynamics and stability correlates of
        global conscious states} (Communications Biology, 2025, strong as
        biology background, medium for behavioural transfer).
  \item \emph{Falling asleep follows a predictable bifurcation dynamic}
        (Nature Neuroscience, 2025, strong for sleep-state transition).
\end{itemize}
\textbf{Reading rule}
\begin{itemize}[nosep]
  \item Use the modern sources to discipline the mechanism language of the
        cases.
  \item Use the classic textbooks to keep the cases teachable and connected to
        durable concepts rather than only recent paper terminology.
  \item Keep frontier consciousness-physics proposals outside the case layer;
        they are too speculative for ordinary behavioural diagnosis.
\end{itemize}
\end{readingbox}

\begin{summarybox}
The case library now does more than illustrate the framework. It functions as
an application-layer bridge between the research spine and practical transfer.
Each case identifies dominant variables, names the relevant path-dependent or
state-transition mechanism, marks the evidence lane clearly, and extracts a
repair pattern that can generalize beyond the anecdote. That is what makes the
chapter feel like an analytical library rather than a set of motivational
stories.
\end{summarybox}
